\documentclass[../main.tex]{subfiles}

\begin{document}

\ifSubfilesClassLoaded{

    %\tableofcontents
    
    %\setcounter{chapter}{}
}{}

\chapter{ Representation of \(  \mathfrak{sl}\left( 2,\mathbb{C}  \right)   \)  }
% 代靖涵 25120222201319

\section{The Modules \(  V_{d}  \) }

\section{Classifying the Irreducible \(  \mathfrak{sl}\left( 2,\mathbb{C}  \right)   \)-Modules }

\begin{lemma}{}{}
    Suppose that \(  V  \) is an \(  \mathfrak{sl}\left( 2,\mathbb{C}  \right)   \)-module and \(  v \in V  \) is an eigenvector of \(  h  \) with eigenvalue \(   \lambda   \).
    \begin{itemize}
        \item Either \(  e\cdot v= 0  \) or \(  e\cdot v  \) is an eigenvector of \(  h  \) with eigenvalue \(   \lambda + 2  \)    .
        \item Either \(  f\cdot v= 0  \) or \(  f\cdot v  \) is an eigenvector of \(  h  \) with eigenvalue \(   \lambda -2  \).    
    \end{itemize}
         
\end{lemma}

\begin{lemma}{}{}
    Let \(  V  \) be a finite-dimensional \(  \mathfrak{sl}\left( 2,\mathbb{C}  \right)   \)-module. Then \(  V  \) contains eigenvector \(  w  \) for \(  h  \) such that \(  e\cdot w= 0  \).      
\end{lemma}

\begin{theorem}{}{}
    If \(  V  \) is finite-dimensional irreducible \(  \mathfrak{sl}\left( 2,\mathbb{C}  \right)   \)-module, then \(  V  \) is isomorphic to one of the \(  V_{d}  \).    
\end{theorem}

\begin{corollary}{}{}
    If \(  V  \) is a finite-dimensional representation of \(  \mathfrak{sl}\left( 2,\mathbb{C}  \right)   \) and \(  w\in V  \) is an  \(  h  \)-eigenvector such that \(  e\cdot w= 0  \), then \(  h\cdot w= dw  \) for some non-negative interger \(  d  \) and the submodule of \(  V  \) generated by \(  w  \) is isomorphic to \(  V_{d}  \).          
\end{corollary}

\section{Weyl's Theorem}

\begin{theorem}{}{}
    Let \(  L  \) be a complex semisimple Lie algebra. Every finite-dimensional representation of \(  L  \) is completely reducible.  
\end{theorem}
\end{document}